\chapter{立体图形与体积}\label{ch:g9:solids}

立体几何从一组基本的\omterm{def:g5:solids:def}{立体图形}开始——\omterm{def:g7:areas:prism}{棱柱}、
\omterm{def:g7:areas:prism}{圆柱}、\omterm{def:g8:solids:def}{棱锥}、\omterm{def:g8:solids:def}{圆锥}、球——并提出两个问题：它们的\omterm{def:g6:measure:volume}{体积}有多大，
切开它们时看到的是怎样的\omterm{prop:g9:solids:sections}{截面}？
本章最后给出一个非常实用的结论：把一个
\omterm{def:g5:solids:def}{立体图形}按比例 $k$ 缩放，\omterm{def:g6:measure:volume}{体积}就变为原来的 $k^3$ 倍。

\section{经典立体图形及其体积}

\begin{theorem}[体积公式]\label{thm:g9:solids:volumes}
记底\omterm{def:g6:measure:area}{面积}为 $B$，高为 $h$（即
底面与对面或\omterm{def:g5:solids:def}{顶点}之间的距离）。
\begin{center}
\renewcommand{\arraystretch}{1.5}
\begin{tabular}{l|l}
\omterm{def:g5:solids:def}{立体图形} & \omterm{def:g6:measure:volume}{体积} \\
\hline
\omterm{def:g7:areas:prism}{棱柱}、\omterm{def:g7:areas:prism}{圆柱} & $V = B \times h$ \\[5pt]
\omterm{def:g8:solids:def}{棱锥}、\omterm{def:g8:solids:def}{圆锥} & $V = \dfrac13\, B \times h$ \\[5pt]
\omterm{def:g3:shapes:shapes}{半径}为 $r$ 的球 & $V = \dfrac43\,\pi r^3$ \\
\end{tabular}
\end{center}
对于\omterm{def:g3:shapes:shapes}{半径}为 $r$ 的\omterm{def:g7:areas:prism}{圆柱}和\omterm{def:g8:solids:def}{圆锥}，底\omterm{def:g6:measure:area}{面积}为 $B = \pi r^2$；
球的表\omterm{def:g6:measure:area}{面积}为 $4\pi r^2$。
\end{theorem}
\admitted

\begin{omfigure}
\begin{tikzpicture}[scale=0.62]
  % cylinder
  \draw[thick] (0,0) ellipse (1 and 0.3);
  \draw[thick] (-1,0) -- (-1,2.6);
  \draw[thick] (1,0) -- (1,2.6);
  \draw[thick] (0,2.6) ellipse (1 and 0.3);
  \draw[dashed] (0,0) -- (0,2.6);
  \node[below, font=\small] at (0,-0.4) {圆柱};
  % cone
  \begin{scope}[xshift=4.2cm]
  \draw[thick] (0,0) ellipse (1 and 0.3);
  \draw[thick] (-1,0) -- (0,2.8);
  \draw[thick] (1,0) -- (0,2.8);
  \draw[dashed] (0,0) -- (0,2.8);
  \node[below, font=\small] at (0,-0.4) {圆锥};
  \end{scope}
  % pyramid
  \begin{scope}[xshift=8.4cm]
  \draw[thick] (-1.1,0) -- (0.9,0) -- (1.5,0.55) -- (-0.5,0.55) -- cycle;
  \draw[thick] (-1.1,0) -- (0.2,2.8);
  \draw[thick] (0.9,0) -- (0.2,2.8);
  \draw[thick] (1.5,0.55) -- (0.2,2.8);
  \draw[dashed] (-0.5,0.55) -- (0.2,2.8);
  \node[below, font=\small] at (0.2,-0.4) {棱锥};
  \end{scope}
  % sphere
  \begin{scope}[xshift=12.8cm]
  \draw[thick] (0,1.3) circle (1.3);
  \draw[dashed] (0,1.3) ellipse (1.3 and 0.4);
  \node[below, font=\small] at (0,-0.4) {球};
  \end{scope}
\end{tikzpicture}

{\small 本章的\omterm{def:g5:solids:def}{立体图形}。每种图形的\omterm{def:g6:measure:volume}{体积}都涉及底\omterm{def:g6:measure:area}{面积}与
高，唯独球例外，只需要\omterm{def:g3:shapes:shapes}{半径}。}
\end{omfigure}

\begin{example}\label{ex:g9:solids:volumes}
一个\omterm{def:g7:areas:prism}{圆柱}形罐头的\omterm{def:g3:shapes:shapes}{半径}为 $4$ cm，高为 $10$ cm：
\[
V = \pi r^2 h = \pi \times 16 \times 10 = 160\pi \approx 503 \text{ cm}^3
\approx 0.5 \text{ L}.
\]
与之底面、高均相同的\omterm{def:g8:solids:def}{圆锥}，\omterm{def:g6:measure:volume}{体积}为其三分之一：
$\frac{160\pi}{3} \approx 168$ cm$^3$。\omterm{def:g3:shapes:shapes}{半径}为 $4$ cm 的球：
$V = \frac43 \pi \times 64 = \frac{256\pi}{3} \approx 268$ cm$^3$。
\end{example}

\begin{example}[逐步求棱锥的体积]\label{ex:g9:solids:pyramid}
一个\omterm{def:g8:solids:def}{棱锥}的底面是边长为 $6$ m 的正方形，高为 $5$ m。
\begin{enumerate}
  \item 底\omterm{def:g6:measure:area}{面积}：$B = 6^2 = 36$ m$^2$。
  \item \omterm{def:g6:measure:volume}{体积}：$V = \frac13 \times 36 \times 5 = 60$ m$^3$。
\end{enumerate}
注意单位：若边长以 cm 为单位、高以 m 为单位，
则须先将其中一个量换算成统一的单位。
\end{example}

\section{平面截立体}

\begin{proposition}[经典立体图形的截面]\label{prop:g9:solids:sections}
用平面截\omterm{def:g5:solids:def}{立体图形}，所得的平面图形称为它的
\emph{截面}\index{截面}：
\begin{itemize}
  \item \omterm{def:g7:areas:prism}{棱柱}或\omterm{def:g7:areas:prism}{圆柱}被\emph{\omterm{def:g6:lines:perp}{平行}于底面}的平面所截，截得
        的图形与底面全等，且在任意高度处都成立；
  \item \omterm{def:g8:solids:def}{棱锥}或\omterm{def:g8:solids:def}{圆锥}被\omterm{def:g6:lines:perp}{平行}于底面的平面所截，截得的图形是
        底面的\emph{缩小图形}：在距\omterm{def:g5:solids:def}{顶点}为 $d$ 处，
        相似比为 $k = \frac{d}{h}$；
  \item \omterm{def:g3:shapes:shapes}{半径}为 $r$ 的球被与球心距离为 $d < r$ 的平面所截，
        得到的\omterm{prop:g9:solids:sections}{截面}是\omterm{def:g3:shapes:shapes}{半径}为 $\sqrt{r^2 - d^2}$ 的圆（由
        勾股定理可得）。
\end{itemize}
\end{proposition}
\admitted

\begin{omfigure}
\begin{tikzpicture}[scale=0.75]
  % cone with section
  \draw[thick] (0,0) ellipse (1.5 and 0.45);
  \draw[thick] (-1.5,0) -- (0,4.0);
  \draw[thick] (1.5,0) -- (0,4.0);
  \draw[omThm, thick, fill=omThm!12] (0,2.4) ellipse (0.6 and 0.18);
  \draw[dashed] (0,0) -- (0,4.0);
  \node[right, font=\small] at (0.1,4.0) {顶点};
  \draw[Stealth-Stealth, thin] (2.1,0) -- (2.1,4.0);
  \node[right, font=\small] at (2.1,2.0) {$h$};
  \draw[Stealth-Stealth, thin, omThm] (0.9,2.4) -- (0.9,4.0);
  \node[right, font=\small, omThm] at (0.92,3.2) {$d$};
  % sphere with section
  \begin{scope}[xshift=7.6cm]
  \draw[thick] (0,2) circle (1.8);
  \draw[omThm, thick, fill=omThm!12] (0,2.9) ellipse (1.56 and 0.4);
  \fill (0,2) circle (1.5pt) node[below, font=\small] {$O$};
  \draw[thin] (0,2) -- (0,2.9) node[midway, left, font=\small] {$d$};
  \draw[thin] (0,2.9) -- (1.56,2.9)
    node[midway, above, font=\small] {$\sqrt{r^2 - d^2}$};
  \draw[thin] (0,2) -- (-1.27,0.72)
    node[midway, below right=-3pt, font=\small] {$r$};
  \end{scope}
\end{tikzpicture}

{\small 左：沿与底面\omterm{def:g6:lines:perp}{平行}、距\omterm{def:g5:solids:def}{顶点} $d$
的平面切割\omterm{def:g8:solids:def}{圆锥}，得到按比例 $k = \frac dh$ 缩放的圆盘。右：一个平面
与球心相距 $d$，截球所得圆的\omterm{def:g3:shapes:shapes}{半径}为
$\sqrt{r^2 - d^2}$。}
\end{omfigure}

\begin{example}\label{ex:g9:solids:section}
一个\omterm{def:g8:solids:def}{圆锥}的底面\omterm{def:g3:shapes:shapes}{半径}为 $6$，高为 $9$。用\omterm{def:g6:lines:perp}{平行}于底面的平面
距\omterm{def:g5:solids:def}{顶点} $3$ 处截得的\omterm{prop:g9:solids:sections}{截面}，是按
$k = \frac39 = \frac13$ 缩放的圆盘，其\omterm{def:g3:shapes:shapes}{半径}为 $6 \times \frac13 = 2$。

一个\omterm{def:g3:shapes:shapes}{半径}为 $5$ 的球被距球心 $3$ 的平面所截，
截得的圆的\omterm{def:g3:shapes:shapes}{半径}为 $\sqrt{25 - 9} = 4$。
\end{example}

\section{缩放立体图形}

\begin{theorem}[缩放对长度、面积与体积的影响]\label{thm:g9:solids:scaling}
\omterm{def:g5:solids:def}{立体图形}按\omterm{def:g9:thales:scaling}{比例因子} $k > 0$ 放大或缩小时：
\begin{itemize}
  \item 所有长度乘以 $k$；
  \item 所有\omterm{def:g6:measure:area}{面积}乘以 $k^2$；
  \item 所有\omterm{def:g6:measure:volume}{体积}乘以 $k^3$。
\end{itemize}
\end{theorem}
\admitted

\begin{example}\label{ex:g9:solids:scaling}
比例为 $\frac{1}{18}$ 的汽车模型：各长度为真实汽车对应长度乘以
$k = \frac{1}{18}$，涂漆表\omterm{def:g6:measure:area}{面积}乘以
$k^2 = \frac{1}{324}$，\omterm{def:g6:measure:volume}{体积}乘以 $k^3 = \frac{1}{5832}$。

将球的\omterm{def:g3:shapes:shapes}{半径}加倍（$k = 2$），其\omterm{def:g6:measure:volume}{体积}变为原来的
$2^3 = 8$ 倍：用公式验证，
$\frac43\pi (2r)^3 = \frac43 \pi \times 8r^3$。
\end{example}

\begin{example}[圆台]\label{ex:g9:solids:truncated}
高为 $9$、底面\omterm{def:g3:shapes:shapes}{半径}为 $6$ 的\omterm{def:g8:solids:def}{圆锥}（\omterm{def:g6:measure:volume}{体积}
$\frac13 \pi \times 36 \times 9 = 108\pi$）在距\omterm{def:g5:solids:def}{顶点} $3$ 处被截开，
\omterm{prop:g9:solids:sections}{截面}以上的小\omterm{def:g8:solids:def}{圆锥}被移去。该小\omterm{def:g8:solids:def}{圆锥}是原\omterm{def:g8:solids:def}{圆锥}按
$k = \frac13$ 缩小所得，因此其\omterm{def:g6:measure:volume}{体积}为
$108\pi \times \left(\frac13\right)^3 = 4\pi$。剩下的\omterm{def:g5:solids:def}{立体图形}（
\emph{圆台}）的\omterm{def:g6:measure:volume}{体积}为 $108\pi - 4\pi = 104\pi \approx 327$。
\end{example}

\section{练习}

\begin{exercise}[$\star$]\label{exo:g9:solids:1}
计算\omterm{def:g6:measure:volume}{体积}：\omterm{def:g5:solids:def}{长方体}的长、宽、高分别为 $4 \times 5
\times 12$；\omterm{def:g7:areas:prism}{圆柱}的\omterm{def:g3:shapes:shapes}{半径}为 $3$、高为 $7$（先给出含
$\pi$ 的精确值，再\omterm{def:g6:decimals:rounding}{四舍五入}到个位）。
\end{exercise}

\begin{exercise}[$\star$]\label{exo:g9:solids:2}
计算底面\omterm{def:g3:shapes:shapes}{半径}为 $5$、高为 $12$ 的\omterm{def:g8:solids:def}{圆锥}的\omterm{def:g6:measure:volume}{体积}，以及
底面为 $8 \times 3$ 的矩形、高为 $10$ 的\omterm{def:g8:solids:def}{棱锥}的\omterm{def:g6:measure:volume}{体积}。
\end{exercise}

\begin{exercise}[$\star$]\label{exo:g9:solids:3}
计算\omterm{def:g3:shapes:shapes}{半径}为 $6$ 的球的\omterm{def:g6:measure:volume}{体积}和表\omterm{def:g6:measure:area}{面积}（精确值
用含 $\pi$ 的形式表示）。
\end{exercise}

\begin{exercise}[$\star$]\label{exo:g9:solids:4}
一个\omterm{def:g3:shapes:shapes}{半径}为 $10$ 的球被与其球心相距 $8$ 的平面所截。
\omterm{prop:g9:solids:sections}{截面}圆的\omterm{def:g3:shapes:shapes}{半径}是多少？
\end{exercise}

\begin{exercise}[$\star\star$]\label{exo:g9:solids:5}
一个\omterm{def:g8:solids:def}{圆锥}的底面\omterm{def:g3:shapes:shapes}{半径}为 $8$，高为 $12$。一个与底面\omterm{def:g6:lines:perp}{平行}的平面
距\omterm{def:g5:solids:def}{顶点} $9$ 处截该\omterm{def:g8:solids:def}{圆锥}。求\omterm{prop:g9:solids:sections}{截面}的\omterm{def:g3:shapes:shapes}{半径}，
再求\omterm{prop:g9:solids:sections}{截面}以上的小\omterm{def:g8:solids:def}{圆锥}的\omterm{def:g6:measure:volume}{体积}。
\end{exercise}

\begin{exercise}[$\star\star$]\label{exo:g9:solids:6}
一个内\omterm{def:g3:shapes:shapes}{半径}为 $3$ cm 的\omterm{def:g7:areas:prism}{圆柱}形玻璃杯，装水至高度
为 $10$ cm。将一个\omterm{def:g3:shapes:shapes}{半径}为 $2$ cm 的小球完全浸入水中，
水面上升了多少？（增加的\omterm{def:g6:measure:volume}{体积}均匀分布在玻璃杯的底\omterm{def:g6:measure:area}{面积}上；
先给出精确的上升高度，再\omterm{def:g6:decimals:rounding}{四舍五入}到毫米。）
\end{exercise}

\begin{exercise}[$\star\star$]\label{exo:g9:solids:7}
一个配方恰好可注满\omterm{def:g3:shapes:shapes}{半径}为 $6$ cm 的球形模具。你只有一个
\omterm{def:g3:shapes:shapes}{半径}为 $3$ cm 的球形模具。可以注满多少个这样的小模具？
（无需显式计算任何一个\omterm{def:g6:measure:volume}{体积}即可作答。）
\end{exercise}

\begin{exercise}[$\star\star$]\label{exo:g9:solids:8}
吉萨大金字塔底面为边长约 $230$ m 的正方形，高约
$147$ m。估算其\omterm{def:g6:measure:volume}{体积}，并以百万
立方米为单位表示。
\end{exercise}

\begin{exercise}[$\star\star\star$]\label{exo:g9:solids:9}
一个\omterm{def:g5:solids:def}{顶点}朝下、底面\omterm{def:g3:shapes:shapes}{半径}为 $6$ cm、高为 $12$ cm 的\omterm{def:g8:solids:def}{圆锥}形漏斗，
装水至其\emph{高度}的\omterm{def:g3:division:half}{一半}（从
\omterm{def:g5:solids:def}{顶点}量起）。
\begin{enumerate}
  \item 被注入的水占漏斗\omterm{def:g6:measure:volume}{体积}的几分之几？
  \item 若改为注入漏斗\emph{\omterm{def:g6:measure:volume}{体积}}\omterm{def:g3:division:half}{一半}的水，
        证明水深 $d$ 满足
        $\left(\frac{d}{12}\right)^3 = \frac12$，并求出 $d$（精确到毫米，
        $\sqrt[3]{0.5} \approx 0.7937$）。
\end{enumerate}
\end{exercise}

\section{问题：阿基米德的墓碑}

\begin{problem}[{周末问题——球恰为其外切圆柱的三分之二
（两次）、1:2:3 的叠放，以及禁止巨人存在的平方--立方定律
}]
\label{pb:g9:solids:1}
阿基米德证明了许多定理，但其中一个让他非常自豪，
以至于他要求将它的\emph{图形}刻在自己的墓碑上：一个球
内切于其最紧的\omterm{def:g7:areas:prism}{圆柱}之中。一个世纪之后，罗马作家
西塞罗在锡拉库萨附近的荆棘丛中搜寻时，认出了
这座墓：“凭球与\omterm{def:g7:areas:prism}{圆柱}”。本题将还原
墓碑所刻内容——双重三分之二的奇迹——然后
依次研究\omterm{def:g6:measure:volume}{体积}与表\omterm{def:g6:measure:area}{面积}（\cref{thm:g9:solids:volumes}、
\cref{thm:g9:solids:scaling}），
进而得到一条支配巨人、蚂蚁与冷却中行星的定律。

\medskip
\noindent\textbf{第一部分——解读墓碑雕刻。}
一个\omterm{def:g3:shapes:shapes}{半径}为 $r$ 的球恰好位于一个\omterm{def:g7:areas:prism}{圆柱}之内：
\omterm{def:g7:areas:prism}{圆柱}与球\omterm{def:g3:shapes:shapes}{半径}相同，高为 $2r$。
\begin{enumerate}
  \item 用 $r$ 表示\omterm{def:g7:areas:prism}{圆柱}的\omterm{def:g6:measure:volume}{体积}。
  \item 求球的\omterm{def:g6:measure:volume}{体积}与
        \omterm{def:g7:areas:prism}{圆柱}\omterm{def:g6:measure:volume}{体积}之比。阿基米德为何会喜欢
        答案中既不含 $\pi$ 也不含 $r$？
  \item 再看表\omterm{def:g6:measure:area}{面积}：将球的表\omterm{def:g6:measure:area}{面积}与
        \omterm{def:g7:areas:prism}{圆柱}的\emph{全部}表\omterm{def:g6:measure:area}{面积}（侧\omterm{def:g6:measure:area}{面积}加上
        两个底面）。再次出现的是什么比值？
  \item 球与\omterm{def:g7:areas:prism}{圆柱}之间的空余部分\omterm{def:g6:measure:volume}{体积}为
        $2\pi r^3 - \frac43\pi r^3$。证明此剩余\omterm{def:g6:measure:volume}{体积}
        恰好等于\emph{两个\omterm{def:g8:solids:def}{圆锥}}的\omterm{def:g6:measure:volume}{体积}，每个\omterm{def:g8:solids:def}{圆锥}
        的\omterm{def:g3:shapes:shapes}{半径}为 $r$、高为 $r$。
  \item 一个\omterm{def:g3:shapes:shapes}{半径}为 $12$ cm 的篮球，装在最紧的
        \omterm{def:g7:areas:prism}{圆柱}形盒子中出售。求两者的\omterm{def:g6:measure:volume}{体积}（精确到
        $100$ cm$^3$），以及盒子中空余部分
        所占的百分比。
\end{enumerate}

\medskip
\noindent\textbf{第二部分——碗、模具与月亮。}
\begin{enumerate}[resume]
  \item 一个半球形碗的\omterm{def:g3:shapes:shapes}{半径}为 $10$ cm。求它的
        \omterm{def:g2:measure:units}{容积}，分别以 cm$^3$ 和升为单位（精确到分升）。
  \item 1:2:3 的叠放：一个\omterm{def:g8:solids:def}{圆锥}、一个半球与一个\omterm{def:g7:areas:prism}{圆柱}，
        \omterm{def:g3:shapes:shapes}{半径}均为 $10$ cm、高也均为 $10$ cm。求三者的
        \omterm{def:g6:measure:volume}{体积}，并验证那个著名的比例
        $1 : 2 : 3$。（阿基米德大概也会欣赏
        这个问题。）
  \item 一个\omterm{def:g3:shapes:shapes}{半径}为 $3$ cm 的巧克力球被熔化后倒入
        \omterm{def:g3:shapes:shapes}{半径}为 $3$ cm 的\omterm{def:g7:areas:prism}{圆柱}形模具中。巧克力的高度
        达到多少？（先给出精确\omterm{def:g9:fractions:fraction}{分数}，再精确到
        毫米。）
  \item 地球的\omterm{def:g3:shapes:shapes}{半径}约为 $6\,371$ km，月球的\omterm{def:g3:shapes:shapes}{半径}约为
        $1\,737$ km。求两者\omterm{def:g3:shapes:shapes}{半径}之比，再
——借助\cref{thm:g9:solids:scaling}——求
        两者的\omterm{def:g6:measure:volume}{体积}之比：按\omterm{def:g6:measure:volume}{体积}计算，一个空心
        地球中能装下多少个月球？
  \item 由同一定理，地球与月球的
        \emph{表\omterm{def:g6:measure:area}{面积}}之比是多少？一般地，气球\omterm{def:g3:shapes:shapes}{半径}
        加倍时，其\omterm{def:g6:measure:volume}{体积}与表\omterm{def:g6:measure:area}{面积}各会如何变化？
        请叙述\emph{平方立方定律}：当图形按比例放大时，
        \omterm{def:g6:measure:volume}{体积}的增长快于表\omterm{def:g6:measure:area}{面积}。
\end{enumerate}

\medskip
\noindent\textbf{第三部分——平方立方定律支配着
世界。}
\begin{enumerate}[resume]
  \item 伽利略反驳巨人存在的论证：设想一个人
        保持比例不变，线性尺度放大 $10$ 倍。重量
        随\omterm{def:g6:measure:volume}{体积}变化：重量增大为原来的多少倍？
        骨骼的横\omterm{prop:g9:solids:sections}{截面}（一个\omterm{def:g6:measure:area}{面积}）增大为
        原来的多少倍？那么，每平方厘米骨骼所承受的\emph{压力}
        增大为原来的多少倍——
        巨人将会怎样？
  \item 同一定律反过来解释蚂蚁的惊人本领：力量
        与肌肉的横\omterm{prop:g9:solids:sections}{截面}（一个\omterm{def:g6:measure:area}{面积}）成正比，重量
        与\omterm{def:g6:measure:volume}{体积}成正比。对一只在各方向都缩小 $100$ 倍的动物，
        重量与力量各缩小到原来的几分之一，
        力量与重量之比则
        提高为原来的多少倍？
  \item 热量由身体的\omterm{def:g6:measure:volume}{体积}产生，经由其表面散失。
        证明：对球体而言，表\omterm{def:g6:measure:area}{面积}与\omterm{def:g6:measure:volume}{体积}之比为
        $\frac{3}{r}$，并计算 $r = 1$ 与
        $r = 2$ 时的值。老鼠与熊哪个散热更快？
        冬天婴儿为何需要戴帽子？
  \item 两个球形橙子的\omterm{def:g3:shapes:shapes}{半径}分别为 $4$ cm 与 $5$ cm，
        价格相同。求二者的\omterm{def:g6:measure:volume}{体积}之比，
        即同样的钱，大橙子能多提供
        多少橙肉？
  \item 总结：请精确描述墓碑上的图形编码了什么
——即第 2 题与第 3 题中的两个“三分之二”
——并用一句话回答西塞罗：一位数学家为何
        胜过他工程天才的一切征服，宁愿选择
        “\omterm{def:g7:areas:prism}{圆柱}中的球”？
\end{enumerate}
\end{problem}
